\clearpage
\fancyhead[L]{\small EFFEMERIDI ASTRONOMICHE.}
\SourceLeaf{561}
\begin{center}\small Göttingische gelehrte Anzeigen. Stück 95. Seite 945..952. 1813 Juni 14.\\[3pt]\rule{.72\linewidth}{.3pt}\end{center}

\emph{Effemeridi astronomiche di Milano per l'anno bisestile} 1812 \emph{calcolate da} \textsc{Francesco Carlini e Carlo Brioschi}. \emph{Con appendice.} 1811. Della reale stamperia. Milano. Der Kalender 124 Seiten, die Zusätze eben so viel, in klein Quart.

\emph{Effemeridi astronomiche di Milano per l'anno} 1813 u. s. f. 1812. Der Kalender 95 Seiten, die Zusätze 136 Seiten in klein Quart.

Wenig bekannt sind bei uns die Mailänder astronomischen Ephemeriden, obgleich dieselben sowohl durch ihre musterhafte Einrichtung, als durch ihre gehaltreichen Zusätze allen ähnlichen, gegenwärtig erscheinenden, Werken den Rang streitig machen. Schade nur, dass sie immer nur kurze Zeit vor Anfang des Jahres, für welches sie bestimmt sind, erscheinen, und daher bei uns gewöhnlich erst im Laufe desselben anlangen. Möchten doch in dieser Hinsicht sich die Herausgeber Hrn. \textsc{Bode} zum Muster nehmen, welcher seine Jahrbücher immer mehr als zwei Jahre vorher erscheinen lässt!

Die Einrichtung des Kalenders ist folgende. Vorausgeschickt sind, ausser dem Schlüssel der Abbreviaturen und einigen Notizen, die sich auf die Zeitrechnung und Kirchenfeste beziehen, eine allgemeine Anzeige der Finsternisse, die scheinbare Schiefe der Ekliptik, und die Nutation der Aequinoctialpunkte in der Länge, beide mit Inbegriff der Solarnutation. Dann folgen die einzelnen Monate, jedem sind 6 Seiten gewidmet. Die erste Seite liefert die Mondsphasen, die Zusammenkunft des Mondes mit Fixsternen, wirkliche Bedeckungen derselben vom Monde, nahe Zusammenkünfte der Planeten mit Fixsternen, und sonst merkwürdige Momente im Planeten- und Sonnenlauf; endlich die Finsternisse der Jupiterstrabanten in mittlerer Zeit. Die zweite Seite zeigt den Abstand der Tage vom Anfang des Jahres, und die Wochentage, mittlere und Sternzeit im wahren Mittage, Sternzeit im mittleren Mittage, Aufgang und Untergang der Sonne. Die dritte Seite enthält die Länge der Sonne, auf Zehntel von Secunden mit grösster Sorgfalt berechnet, die gerade Aufsteigung der Sonne in Bogen, ihre Declination, und die Logarithmen ihres Abstandes von der Erde mit sechs Decimalen, letztern für alle einzelne Tage, welches ein sehr grosser, diesen Ephemeriden eigenthümlicher, Vorzug ist. Die vierte und fünfte Seite ist dem Monde gewidmet, und liefert Länge, Breite, Horizontalparallaxe und Durchmesser desselben sowohl für Mittag als für Mitternacht, die Culminationszeit, Aufgang und Untergang (nur in Zeitminuten), und die Declination für den Durchgang durch den Meridian in Bogenminuten. Auf der sechsten Seite endlich sind für die einzelnen Tage des Monats die Configurationen der Jupiterstrabanten abgebildet. Dann folgen Halbmesser der Sonne, Culminationsdauer derselben und Länge des Mondsknoten von 6 zu 6 Tagen durch das ganze Jahr fortlaufend, und hierauf die Bewegungen der einzelnen zehn Planeten: dass letztere nicht stückweise nach den einzelnen Monaten, sondern in Einer Uebersicht für das ganze Jahr zusammengestellt sind, finden wir sehr zweckmässig. Mercur, Venus, Mars füllen jeder zwei Seiten, da die Angaben für dieselben von 6 zu 6 Tagen durch das ganze Jahr fortlaufen; Ceres, Pallas, Juno, Vesta, Jupiter, Saturn und Uranus nur halb so viel, indem die Oerter der vier erstern nur für die Zeit ihrer Sichtbarkeit, die der letztern nur von 12 zu 12 Tagen angesetzt sind. Bei jedem einzelnen Planeten wird (auf Minuten) angegeben Länge, Breite, gerade Aufsteigung (in Zeitminuten), Declination, Aufgang, Culmination und Untergang. Dürften wir uns hierbei noch einen Wunsch erlauben, so wäre es, dass die geraden Aufsteigungen bis auf Bogenminuten genau angesetzt sein möchten (welches besonders bei den neuen Planeten von Wichtigkeit ist), und dass ausserdem den Abständen der Planeten von der Erde eine eigene Columne gewidmet wäre, die besonders für die Berechnung der Aberration, so wie für die \SourceLeaf{562}Parallaxe, scheinbaren Durchmesser und andere Zwecke nützlich sein würden, und wofür, um Raum zu gewinnen, der Aufgang und der Untergang der Planeten, je nachdem sie in den Abend- oder Frühstunden culminiren, wegfallen könnte, so wie die Columne für die Länge etwas schmaler ausfallen würde, wenn diese nicht in Zeichen, Graden und Minuten, sondern nur in Graden und Minuten abgedruckt wäre. Noch bemerken wir, dass die geraden Aufsteigungen der Pallas im Jahre 1813, sowie die davon abhängenden Culminations-Zeiten, durch einen Rechnungsfehler alle falsch sind, und dass für jene die Complemente zu 42 Stunden genommen werden müssen. --- Zuletzt sind noch die Bedeckungen der Fixsterne nach den Rechnungen der Florenzer Astronomen mitgetheilt, eben so, wie sie alljährlich auch in der Monatl. Correspondenz abgedruckt werden. Ausser diesen stehenden Artikeln sind, in dem Jahrgange für 1812, dem Kalender noch beigefügt: Ein Verzeichniss aller in Mailand sichtbaren Sterne über der fünften Grösse, nach \textsc{Piazzi} auf 181c reducirt (zusammen 596), und die \textsc{Gauss}ischen Tafeln für Aberration und Nutation, nebst einer kleinen Tafel für die Solarnutation.

Unter den \emph{Zusätzen} zum Jahrgange 1812 nehmen die Beobachtungen von Zenithdistanzen der Sonne und Fixsterne im Meridian, mit einem neuen Wiederholungskreise, von \textsc{Barnabas Oriani}, den ersten Platz ein. Voraus geschickt ist eine lehrreiche Beschreibung des unvergleichlichen Instruments, eines dreifüssigen \textsc{Reichenbach}'schen Kreises mit stehender Säule (Preis 3000 Gulden). Die zweijährigen, mit diesem Instrumente angestellten und in den beiden Jahrgängen 1812 und 1813 abgedruckten, Beobachtungen enthalten für die Sterndeclinationen, für die Theorie der Bewegung der Sonne, und für die Theorie der astronomischen Strahlenbrechung, einen Schatz von Erfahrungen, den wir bald ganz so, wie er es verdient, benutzt wünschen. --- Beobachtungen zur Bestimmung der Solstitien und der Schiefe der Ekliptik in den Jahren 1810 und 1811, von \textsc{Angelo Cesaris}. Während \textsc{Oriani} das Wintersolstitium von 1810, und das Sommersolstitium von 1811, mit dem dreifüssigen \textsc{Reichenbach}'schen Kreise beobachtete, bediente sich \textsc{Cesaris} zu demselben Zwecke des achtfüssigen Mauerquadranten. So wie die beiderseitigen Resultate hier mitgetheilt sind, stimmen sie innerhalb einiger Zehntel der Secunde überein; wir hätten nur dabei eine Erklärung darüber gewünscht, auf welchem Wege der Collimationsfehler des Quadranten ausgemittelt worden ist, welchen \textsc{Cesaris} in beiden Solstitien $= -1''5$ annimmt. Auch deuten diese Beobachtungen auf keine Verschiedenheit der Schiefe in dem Winter- und Sommersolstitium hin, welche einige Astronomen gefunden haben wollen, da die von \textsc{Cesaris} unter Voraussetzung von einerlei Schiefe abgeleitete Polhöhe $45^\circ 28' 0''20$ sehr nahe mit der auf anderem Wege gefundenen übereinstimmt. --- Ueber den Grad der Convergenz der verschiedenen Reihen, welche die Ungleichheiten der Mondslänge darstellen, von \textsc{Franz Carlini}. Ein ungemein schätzbarer Aufsatz! Es werden hier drei Ausdrücke für die Mondsungleichheiten zusammengestellt; der erste nach \textsc{Tobias Mayer}'s Form, wonach die \textsc{Bürg}'schen Tafeln berechnet sind; der zweite, aus dem ersten von \textsc{Laplace} abgeleitet, in der Form, wie \textsc{Laplace}'s Theorie diese Ungleichheiten gegeben hat, d. i. als Functionen der wahren Mondslänge; der dritte, von \textsc{Carlini} aus dem ersten berechnete, in der Form, in welche schon \textsc{Lambert} und \textsc{Schulze} die \textsc{Mayer}'schen und \textsc{Mason}'schen Gleichungen gebracht hatten, d. i. blos als Functionen der mittlern Bewegungen. Wir würden diese letztere Form allen andern vorziehen, wenngleich der Ausdruck etwas langsamer convergirt, als die beiden ersten. Es wäre zu wünschen, dass \emph{alle} Coëfficienten desselben unmittelbar aus einigen Tausend Beobachtungen abgeleitet würden, was freilich nur die Astronomen mit Vortheil ausführen könnten, denen die Vorarbeiten \textsc{Bürg}'s oder \textsc{Burckhardt}'s dabei zu Gebote ständen: an 100 aus Sternbedeckungen abgeleiteten Mondsörtern hat \textsc{Carlini} die Prüfung selbst vorgenommen und durchaus gute Uebereinstimmung gefunden. Merkwürdig ist die hier von \textsc{Carlini} gefundene Bestätigung einer neuen, schon früher von \textsc{Burckhardt} aufgestellten, aber jetzt nicht in dessen neuen Tafeln aufgenommenen, Gleichung, die sich zuerst in der Gestalt einer Ungleichheit der Excen\SourceLeaf{563}tricität mit einer langen Periode ankündigt. --- Ueber den Einfluss der Aenderungen der Temperatur auf die Bewegungen des Pendels, von \textsc{Carl Brioschi}. Eine Untersuchung, die der Verfasser mit einem grössern Aufwande von Kunst durchgeführt hat, als sie bedurft hätte. Im Wesentlichen besteht das Resultat doch nur darin, dass in einem gewissen Zeitintervall ein Pendel, dessen Länge bei veränderlicher Temperatur eine veränderliche Länge hat, gerade eben so viele Schwingungen macht, als es mit constanter Länge bei dem mittlern Thermometerstande gemacht haben würde. --- Zuletzt noch einige in Mailand und Rom von \textsc{Carlini} und \textsc{Oriani} in den Jahren 1808 und 1810 beobachtete Sternbedeckungen.

Die Zusätze zu dem Jahrgange 1813 fangen an mit der schon oben erwähnten Fortsetzung der mit dem dreifüssigen \textsc{Reichenbach}'schen Kreise beobachteten Zenithdistanzen. --- Ueber das periodische Schwanken der Gebäude, von \textsc{Angelo Cesaris}. Man weiss längst, wie wenig man sich auf die unverrückte Lage solcher astronomischen Instrumente, die eine feste Aufstellung erfordern, verlassen könne, wenn sie zumal in beträchtlicher Erhöhung über der Erde mit den Aussenmauern des Gebäudes in Verbindung stehen, auf welche Temperatur und Feuchtigkeit ihre unmittelbaren Einwirkungen äussern, und wie wichtig es daher bei Anlegung einer neuen Sternwarte sei, solchen Instrumenten zu ebener Erde eine vollkommen feste, von den äussern Mauern ganz unabhängige, Basis zu geben. Die übrigens mit so vortrefflichen Instrumenten ausgerüstete Sternwarte Brera in Mailand hat diesen Vorzug nicht, und die Schwankungen des Gebäudes zeigen sich an dem Mauerquadranten und Passageninstrumente sehr bestimmt und stark. Wenn das in einer Entfernung von 3000 Toisen errichtete Meridianzeichen von dem Verticalfaden des Passageninstruments Morgens vor Sonnenaufgang berührt wird, so entfernt sich späterhin dieser immer mehr östlich von jenem bis nach Mittag, wo der Abstand an heitern Wintertagen auf 5 bis 6 Secunden, an heitern Sonnentagen aber auf 30 Secunden geht. Unter übrigens gleichen Umständen ist die Veränderung bei bedecktem Himmel und sich fast gleich bleibender Temperatur am geringsten oder ganz unmerklich, hingegen am grössten bei Sonnenschein und starker Temperaturveränderung. --- Ueber die Formeln für die Parallaxe und Breite des Mondes, von \textsc{Franz Carlini}: eine Fortsetzung des Aufsatzes im vorhergehenden Bande. Die Formel für die Parallaxe wird auf ähnliche Art, wie die für die Ungleichheiten der Länge, in eine andere verwandelt, die blos von mittlern Bewegungen abhängt; einen ähnlichen Ausdruck hat \textsc{Carlini} für den Logarithmen der Parallaxe entwickelt. Hingegen hat er den Ausdruck der Breite nicht auch auf dieselbe Art umgeformt, weil diese gar zu ungeschmeidig und langsam convergirend ausgefallen sein würde, und \textsc{Carlini} begnügt sich damit, ein sehr einfaches Verfahren anzugeben, wie Tafeln bequem für die \textsc{Laplace}'sche Form (in Functionen der wahren Länge) eingerichtet werden könnten. Aus mehreren Gründen hätten wir doch gewünscht, dass Carlini auch bei diesem Element, \emph{nur mit Ausschluss der beiden ersten Glieder} 
\[
18540''25\sin\delta + 12''56\sin 3\delta
\]
jene Verwandlung vorgenommen hätte; jene Unbequemlichkeit wäre alsdann weggefallen, und die Analogie zwischen den Formeln für die Ungleichheiten des Mondes und denen für die Ungleichheiten der Planeten wäre dadurch vollständig geworden. Inzwischen kann diese Verwandlung leicht nachgeholt werden. --- Opposition des Saturn im Jahre 1811, beobachtet von \textsc{Carl Brioschi}. --- Auszug aus den meteorologischen Beobachtungen auf der Mailänder Sternwarte im Jahre 1809, von \textsc{Angelo Cesaris}.
